(1) 槽可旋转且无外部方位力矩，珠受到的槽法向约束力并不在外部系统中生成净力矩；把自由旋转槽与珠视为整体可得 \(L_z=mr^2\dot\phi=mr_0^2\omega_0\)。因此 \(\dot\phi=\omega_0/(1+\eta t)^2\)，积分给 \(\phi(t)-\phi(0)=\omega_0 t/(1+\eta t)\)。若固定槽方位不允许自由转动，则会有外部轴力矩，本结论不成立；本题特意规定槽可自由转动。

(2) 质点在极坐标的切向加速度是 \(r\ddot\phi+2\dot r\dot\phi=0\)，所以无须外部切向力来改变角速度；径向加速度 \(a_r=\ddot r-r\dot\phi^2=-r\dot\phi^2\)。执行器施加径向合力 \(F_r=-mr\dot\phi^2=-mr_0\omega_0^2/(1+\eta t)^3\)，为使半径持续增大必须做负功。假设槽无惯性，切向约束力恒零；若槽有惯性则要联合其角动量重新建模。

(3) 动能 \(K=m[\dot r^2+r^2\dot\phi^2]/2=(m/2)[r_0^2\eta^2+r_0^2\omega_0^2/(1+\eta t)^2]\)。径向强制速度恒定，故 \(\Delta K=(mr_0^2\omega_0^2/2)[(1+\eta t)^{-2}-1]\)。外驱功 \(W=\int_0^tF_r\dot r\,dt'=-mr_0^2\eta\omega_0^2\int_0^t(1+\eta t')^{-3}dt'\)，恰等于上述动能变化。负号表示向外伸长时，角向转动动能由执行器回收。

(4) 当 \(t\to\infty\) 时 \(\dot\phi\sim\omega_0/(\eta t)^2\)，角向速度 \(r\dot\phi\sim r_0\omega_0/(\eta t)\)，角动能趋零，但固定速率径向动能仍非零。若执行器反向以无损、准静态的方式收回，径向力方向与速度可能相同而对珠做正功，将已存储的机械能返还；有粘滞摩擦或控制器内部耗散时不能保证完整回收。角动量守恒不等于机械能守恒，因为时变径向约束显式作功。
